\section{Results}
\label{sec:chap_slam_results}

\todo{explain the odometry and icp for ground truth here.}
\begin{itemize}
    \item Note that for the experiments, we used a \textit{C++} implementation developed by Bastian Steder.
    \item We tried with the rotation invariance parameters but result were less conclusive
    \item Confusion matrix: The dark areas that are not close to the main diagonal mark loop closures.
\end{itemize}

\subsection{Structured vs Forest}
\label{ssec:chap_slam_struct_vs_forest}

\begin{figure}[htpb]
    \centering
    \includegraphics[width=0.2\linewidth]{img/todo.jpg}
    \caption{Confusion matrices structured (left) vs forest (right). Notice that the diagonals are wider for the structured environment}
    \label{fig:matrices_struct_vs_forest}
\end{figure}

\begin{figure}[htpb]
    \centering
    \includegraphics[width=0.2\linewidth]{img/todo.jpg}
    \caption{Nice to have: recall vs dist}
    \label{fig:recall_dist}
\end{figure}

\begin{figure}[htpb]
    \centering
    \includegraphics[width=0.2\linewidth]{img/todo.jpg}
    \caption{Nice to have: matched keypoints (something like an histogram of dist, height or avg dist per scan)}
    \label{fig:match_avg_dist}
\end{figure}

\subsection{SICK vs Velodyne}
\label{ssec:chap_slam_sick_vs_velodyne}
